\textbackslash documentclass{[}12pt,a4paper{]}\{article\}
\textbackslash usepackage\{amsmath, amssymb\}
\textbackslash usepackage\{geometry\}
\textbackslash usepackage\{graphicx\}
\textbackslash usepackage\{hyperref\}
\textbackslash usepackage\{setspace\}
\textbackslash geometry\{margin=1in\} \textbackslash setstretch\{1.2\}

\textbackslash title\{CSE400 -\/- Fundamentals of Probability in
Computing \textbackslash\textbackslash{} Lecture L11-\/-L12:
Transformation of Random Variables\} \textbackslash author\{Instructor:
Dhaval Patel, PhD\} \textbackslash date\{\}

\textbackslash begin\{document\} \textbackslash maketitle

\textbackslash section*\{Slide 1: Introduction to Transformation of
Random Variables\}

\textbackslash subsection*\{Main Idea\}

In probability, many times: \textbackslash begin\{itemize\}
\textbackslash item We already know a random variable \$X\$.
\textbackslash item We know its PDF \$f\_X(x)\$ and CDF \$F\_X(x)\$.
\textbackslash item But we are interested in a new variable made from
\$X\$ (or from \$X\$ and \$Y\$). \textbackslash end\{itemize\}

Examples: \textbackslash{[} Y = g(X) \textbackslash{]} \textbackslash{[}
Z = g(X,Y), \textbackslash quad \textbackslash text\{such as \} Z = X+Y,
\textbackslash; Z = X-Y, \textbackslash; Z =
\textbackslash frac\{X\}\{Y\}, \textbackslash; Z = X\textbackslash cdot
Y \textbackslash{]}

\textbackslash subsection*\{Goal of This Lecture\}

If \$X\$ (or \$X,Y\$) is known, how do we find: \textbackslash{[}
F\_Y(y), \textbackslash; f\_Y(y) \textbackslash quad
\textbackslash text\{or\} \textbackslash quad F\_Z(z), \textbackslash;
f\_Z(z)? \textbackslash{]} This process is called
\textbackslash textbf\{Transformation of Random Variables\}.

\textbackslash section*\{Slide 2: Outline of the Lecture\}

\textbackslash begin\{itemize\} \textbackslash item Transformation of
Random Variables (single variable case) \textbackslash item Function of
Two Random Variables (joint transformations) \textbackslash item
Illustrative Example: \$Z = X + Y\$ (full derivation)
\textbackslash end\{itemize\}

The function \$g(\textbackslash cdot)\$ can be monotonically increasing
or decreasing, and limits and signs matter while transforming.

\textbackslash section*\{Slide 3: Case 1 -\/-\/- Transformation of One
Random Variable\}

Let \$X\$ be a continuous random variable with known PDF \$f\_X(x)\$ and
CDF \$F\_X(x)\$. Define: \textbackslash{[} Y = g(X) \textbackslash{]} We
want to find \$F\_Y(y)\$ and \$f\_Y(y)\$.

\textbackslash subsection*\{Step 1: CDF Method\}

By definition: \textbackslash{[} F\_Y(y) = P(Y \textbackslash le y) =
P(g(X) \textbackslash le y) \textbackslash{]}

If \$g(\textbackslash cdot)\$ is monotonically increasing:
\textbackslash{[} g(X) \textbackslash le y \textbackslash iff X
\textbackslash le g\^{}\{-1\}(y) \textbackslash{]} So: \textbackslash{[}
F\_Y(y) = P(X \textbackslash le g\^{}\{-1\}(y)) = F\_X(g\^{}\{-1\}(y))
\textbackslash{]}

\textbackslash subsection*\{Step 2: Differentiate to Get the PDF\}

\textbackslash{[} f\_Y(y) = \textbackslash frac\{d\}\{dy\} F\_Y(y) =
\textbackslash frac\{d\}\{dy\} F\_X(g\^{}\{-1\}(y)) \textbackslash{]}

Using chain rule: \textbackslash{[} f\_Y(y) = f\_X(g\^{}\{-1\}(y))
\textbackslash cdot \textbackslash frac\{d\}\{dy\}
\textbackslash left{[}g\^{}\{-1\}(y)\textbackslash right{]}
\textbackslash{]}

Taking absolute value: \textbackslash{[} \textbackslash boxed\{ f\_Y(y)
= f\_X(x) \textbackslash left\textbar{} \textbackslash frac\{dx\}\{dy\}
\textbackslash right\textbar{} \textbackslash Bigg\textbar\_\{x =
g\^{}\{-1\}(y)\} \} \textbackslash{]}

\textbackslash section*\{Slide 4: If \$g(x)\$ is Decreasing\}

If \$g(\textbackslash cdot)\$ is monotonically decreasing:
\textbackslash{[} F\_Y(y) = P(g(X) \textbackslash le y) = P(X
\textbackslash ge g\^{}\{-1\}(y)) = 1 - F\_X(g\^{}\{-1\}(y))
\textbackslash{]}

Differentiating: \textbackslash{[} f\_Y(y) = - f\_X(g\^{}\{-1\}(y))
\textbackslash cdot \textbackslash frac\{d\}\{dy\}{[}g\^{}\{-1\}(y){]}
\textbackslash{]}

Taking absolute value again gives: \textbackslash{[}
\textbackslash boxed\{ f\_Y(y) = f\_X(g\^{}\{-1\}(y))
\textbackslash left\textbar{} \textbackslash frac\{dx\}\{dy\}
\textbackslash right\textbar{} \} \textbackslash{]}

So the same formula works for both increasing and decreasing cases.

\textbackslash section*\{Slide 5: Example 1 -\/-\/- \$Y =
\textbackslash sin\textbackslash left(\textbackslash frac\{\textbackslash pi
X\}\{2\}\textbackslash right)\$, \$X \textbackslash sim U(-1,1)\$\}

Given: \textbackslash{[} f\_X(x) = \textbackslash frac\{1\}\{2\},
\textbackslash quad -1 \textless{} x \textless{} 1 \textbackslash{]}

\textbackslash subsection*\{Step 1: Find the Inverse\}

\textbackslash{[} y =
\textbackslash sin\textbackslash left(\textbackslash frac\{\textbackslash pi
x\}\{2\}\textbackslash right) \textbackslash Rightarrow x =
\textbackslash frac\{2\}\{\textbackslash pi\}
\textbackslash sin\^{}\{-1\}(y) \textbackslash{]}

\textbackslash subsection*\{Step 2: Compute
\$\textbackslash frac\{dx\}\{dy\}\$\}

\textbackslash{[} \textbackslash frac\{dx\}\{dy\} =
\textbackslash frac\{2\}\{\textbackslash pi\} \textbackslash cdot
\textbackslash frac\{1\}\{\textbackslash sqrt\{1-y\^{}2\}\}
\textbackslash{]}

\textbackslash subsection*\{Step 3: Apply the Formula\}

\textbackslash{[} f\_Y(y) = f\_X(x) \textbackslash left\textbar{}
\textbackslash frac\{dx\}\{dy\} \textbackslash right\textbar{} =
\textbackslash frac\{1\}\{2\} \textbackslash cdot
\textbackslash frac\{2\}\{\textbackslash pi\} \textbackslash cdot
\textbackslash frac\{1\}\{\textbackslash sqrt\{1-y\^{}2\}\} =
\textbackslash frac\{1\}\{\textbackslash pi
\textbackslash sqrt\{1-y\^{}2\}\} \textbackslash{]}

Domain: \$-1 \textless{} y \textless{} 1\$.

\textbackslash{[} \textbackslash boxed\{ f\_Y(y) =
\textbackslash frac\{1\}\{\textbackslash pi
\textbackslash sqrt\{1-y\^{}2\}\}, \textbackslash quad -1 \textless{} y
\textless{} 1 \} \textbackslash{]}

\textbackslash section*\{Slide 6: Case 2 -\/-\/- Function of Two Random
Variables (\$Z = X + Y\$)\}

We want to find: \textbackslash begin\{itemize\} \textbackslash item (i)
PDF of \$Z\$, \$f\_Z(z)\$ \textbackslash item (ii) If \$X\$ and \$Y\$
are independent \textbackslash item (iii) If \$X \textbackslash sim
N(0,1)\$, \$Y \textbackslash sim N(0,1)\$, show \$Z \textbackslash sim
N(0,2)\$ \textbackslash item (iv) If \$X,Y\$ are exponential with
parameter \$\textbackslash lambda\$, find \$f\_Z(z)\$
\textbackslash end\{itemize\}

\textbackslash section*\{Slide 7: CDF of \$Z = X + Y\$\}

\textbackslash{[} F\_Z(z) = P(Z \textbackslash le z) = P(X+Y
\textbackslash le z) \textbackslash{]}

This corresponds to the region: \textbackslash{[} x + y
\textbackslash le z \textbackslash{]}

So: \textbackslash{[} F\_Z(z) = \textbackslash iint\_\{x+y
\textbackslash le z\} f\_\{X,Y\}(x,y) \textbackslash, dx \textbackslash,
dy \textbackslash{]}

\textbackslash section*\{Slide 8: Leibniz Rule\}

If: \textbackslash{[} G(x) = \textbackslash int\_\{a(x)\}\^{}\{b(x)\}
h(x,y) \textbackslash, dy \textbackslash{]} Then: \textbackslash{[}
\textbackslash frac\{dG\}\{dx\} = h(x,b(x)) b\textquotesingle(x) -
h(x,a(x)) a\textquotesingle(x) +
\textbackslash int\_\{a(x)\}\^{}\{b(x)\}
\textbackslash frac\{\textbackslash partial h\}\{\textbackslash partial
x\} dy \textbackslash{]}

\textbackslash section*\{Slide 9: Deriving the PDF of \$Z\$\}

\textbackslash{[} F\_Z(z) = \textbackslash iint\_\{x+y \textbackslash le
z\} f\_\{X,Y\}(x,y) \textbackslash, dx \textbackslash, dy
\textbackslash{]}

Differentiating: \textbackslash{[} \textbackslash boxed\{ f\_Z(z) =
\textbackslash int\_\{-\textbackslash infty\}\^{}\{\textbackslash infty\}
f\_\{X,Y\}(x, z-x) \textbackslash, dx \} \textbackslash{]}

\textbackslash section*\{Slide 10: Independent Case (Convolution)\}

If \$X\$ and \$Y\$ are independent: \textbackslash{[} f\_\{X,Y\}(x,y) =
f\_X(x) f\_Y(y) \textbackslash{]}

So: \textbackslash{[} \textbackslash boxed\{ f\_Z(z) =
\textbackslash int\_\{-\textbackslash infty\}\^{}\{\textbackslash infty\}
f\_X(x) f\_Y(z-x) \textbackslash, dx \} \textbackslash{]} This is called
the \textbackslash textbf\{convolution integral\}.

\textbackslash section*\{Slides 10-\/-11: Sum of Two Standard
Gaussians\}

Given: \textbackslash{[} X \textbackslash sim N(0,1),
\textbackslash quad Y \textbackslash sim N(0,1) \textbackslash{]}
\textbackslash{[} f\_X(x) =
\textbackslash frac\{1\}\{\textbackslash sqrt\{2\textbackslash pi\}\}
e\^{}\{-x\^{}2/2\}, \textbackslash quad f\_Y(y) =
\textbackslash frac\{1\}\{\textbackslash sqrt\{2\textbackslash pi\}\}
e\^{}\{-y\^{}2/2\} \textbackslash{]}

\textbackslash{[} f\_Z(z) =
\textbackslash int\_\{-\textbackslash infty\}\^{}\{\textbackslash infty\}
\textbackslash frac\{1\}\{\textbackslash sqrt\{2\textbackslash pi\}\}
e\^{}\{-x\^{}2/2\} \textbackslash cdot
\textbackslash frac\{1\}\{\textbackslash sqrt\{2\textbackslash pi\}\}
e\^{}\{-(z-x)\^{}2/2\} \textbackslash, dx \textbackslash{]}

After simplifying: \textbackslash{[} f\_Z(z) =
\textbackslash frac\{1\}\{\textbackslash sqrt\{4\textbackslash pi\}\}
e\^{}\{-z\^{}2/4\} \textbackslash{]}

So: \textbackslash{[} \textbackslash boxed\{ Z \textbackslash sim N(0,2)
\} \textbackslash{]}

\textbackslash section*\{Slide 12: \$X\$ and \$Y\$ are
Exponential(\$\textbackslash lambda\$)\}

\textbackslash{[} f\_X(x) = \textbackslash lambda
e\^{}\{-\textbackslash lambda x\}, \textbackslash; x \textbackslash ge
0, \textbackslash quad f\_Y(y) = \textbackslash lambda
e\^{}\{-\textbackslash lambda y\}, \textbackslash; y \textbackslash ge 0
\textbackslash{]}

\textbackslash{[} f\_Z(z) = \textbackslash int\_0\^{}z
\textbackslash lambda e\^{}\{-\textbackslash lambda x\}
\textbackslash cdot \textbackslash lambda
e\^{}\{-\textbackslash lambda(z-x)\} \textbackslash, dx =
\textbackslash lambda\^{}2 e\^{}\{-\textbackslash lambda z\}
\textbackslash int\_0\^{}z dx = \textbackslash lambda\^{}2 z
e\^{}\{-\textbackslash lambda z\} \textbackslash{]}

\textbackslash{[} \textbackslash boxed\{ f\_Z(z) =
\textbackslash lambda\^{}2 z e\^{}\{-\textbackslash lambda z\},
\textbackslash quad z \textbackslash ge 0 \} \textbackslash{]} This is a
Gamma distribution with shape \$2\$.

\textbackslash section*\{Slide 13: Example -\/-\/- \$Z = X - Y\$\}

\textbackslash{[} F\_Z(z) = P(X - Y \textbackslash le z) = P(Y
\textbackslash ge X - z) \textbackslash{]} Integrate the joint PDF over
this region and then differentiate: \textbackslash{[} f\_Z(z) =
\textbackslash frac\{d\}\{dz\} F\_Z(z) \textbackslash{]} Procedure is
similar to \$Z = X + Y\$.

\textbackslash section*\{Slide 14: Examples -\/-\/- \$Z = X/Y\$ and \$R
= \textbackslash sqrt\{X\^{}2 + Y\^{}2\}\$\}

\textbackslash subsection*\{Case 1: \$Z = X/Y\$\}

\textbackslash{[} F\_Z(z) =
P\textbackslash left(\textbackslash frac\{X\}\{Y\} \textbackslash le
z\textbackslash right) \textbackslash{]} If \$Y\textgreater0\$,
inequality stays same. If \$Y\textless0\$, inequality reverses. So
integration is done piecewise over regions.

\textbackslash subsection*\{Case 2: \$R = \textbackslash sqrt\{X\^{}2 +
Y\^{}2\}\$\}

\textbackslash{[} F\_R(r) = P(\textbackslash sqrt\{X\^{}2 + Y\^{}2\}
\textbackslash le r) = P(X\^{}2 + Y\^{}2 \textbackslash le r\^{}2)
\textbackslash{]} This is probability inside a circle of radius \$r\$:
\textbackslash{[} F\_R(r) = \textbackslash iint\_\{x\^{}2 + y\^{}2
\textbackslash le r\^{}2\} f\_\{X,Y\}(x,y) \textbackslash, dx
\textbackslash, dy \textbackslash{]} \textbackslash{[} f\_R(r) =
\textbackslash frac\{d\}\{dr\} F\_R(r) \textbackslash{]}

\textbackslash section*\{Final Summary\}

\textbackslash begin\{itemize\} \textbackslash item One-variable
transformation: \textbackslash{[} \textbackslash boxed\{ f\_Y(y) =
f\_X(g\^{}\{-1\}(y)) \textbackslash left\textbar{}
\textbackslash frac\{dx\}\{dy\} \textbackslash right\textbar{} \}
\textbackslash{]} \textbackslash item Two-variable transformation:
\textbackslash{[} F\_Z(z) =
\textbackslash iint\_\{g(x,y)\textbackslash le z\} f\_\{X,Y\}(x,y)
\textbackslash, dx \textbackslash, dy, \textbackslash quad f\_Z(z) =
\textbackslash frac\{d\}\{dz\} F\_Z(z) \textbackslash{]}
\textbackslash item Sum of RVs: \textbackslash{[} f\_Z(z) =
\textbackslash int\_\{-\textbackslash infty\}\^{}\{\textbackslash infty\}
f\_X(x) f\_Y(z-x) \textbackslash, dx \textbackslash{]}
\textbackslash item Leibniz Rule: \textbackslash{[}
\textbackslash frac\{d\}\{dx\} \textbackslash int\_\{a(x)\}\^{}\{b(x)\}
h(x,y) dy = h(x,b(x)) b\textquotesingle(x) - h(x,a(x))
a\textquotesingle(x) + \textbackslash int\_\{a(x)\}\^{}\{b(x)\}
\textbackslash frac\{\textbackslash partial h\}\{\textbackslash partial
x\} dy \textbackslash{]} \textbackslash end\{itemize\}

\textbackslash end\{document\}
